% Residencia pública de la genealogía anterior a la publicación parangular.
% Este bloque no sustituye ningún propietario: fija el orden de lectura de
% las dos ramas ya demostradas y de su confluencia tipada.

\section{Bifurcation généalogique de l'échelle d'action et de sa composante adimensionnelle}
\label{sec:sucesor-102-c14-bifurcacion-accion}

La phase intrinsèque \(\alpha_{\mathrm{HMT}}\), déjà obtenue par la
clôture dodécaphasique réversible, engendre deux branches internes indépendantes.
La première détermine la composante adimensionnelle de l'action postérieure au
retour ; la seconde détermine sa décennie absolue. Aucune d'elles ne présuppose
le contre-angle.

Dans la branche régionale, le caractère de degré \(5\) est
\begin{equation}
 H_5(\alpha_{\mathrm{HMT}},\varphi)
 =\frac12\sqrt{\frac{10^3\alpha_{\mathrm{HMT}}}{\varphi}}
  -\alpha_{\mathrm{HMT}}
  +\frac9{16}\alpha_{\mathrm{HMT}}^2
  -\frac59\alpha_{\mathrm{HMT}}^3
  +\frac7{48}\alpha_{\mathrm{HMT}}^4
  -\frac1{54}\alpha_{\mathrm{HMT}}^5 .
 \label{eq:sucesor-102-c14-h5-preangular}
\end{equation}
La complétion \(10^3=729+243+27+1\) publie la phase dans la coordonnée
analytique
\begin{equation}
 x_A=\frac{50\pi}{9}\alpha_{\mathrm{HMT}},
 \qquad
 D(\alpha_{\mathrm{HMT}})=e^{-2x_A}
 =\exp\!\left(-\frac{100\pi}{9}\alpha_{\mathrm{HMT}}\right).
 \label{eq:sucesor-102-c14-determinante-regional}
\end{equation}
Avec le sélecteur régional \(\rho_\pi=169\), la récurrence convergente
\begin{equation}
 \mathscr C_\pi(\alpha)
 =169\alpha^6+
   \frac{D(\alpha)\alpha^7}{1-D(\alpha)\alpha}
 \label{eq:sucesor-102-c14-recurrencia-regional}
\end{equation}
détermine directement
\begin{equation}
 \boxed{\eta_{\mathrm{ret}}
 =H_5-\frac{90}{\pi}\mathscr C_\pi(\alpha_{\mathrm{HMT}})>0.}
 \label{eq:sucesor-102-c14-eta-directa}
\end{equation}
Cette équation ne contient pas \(C^*\) : elle construit la coordonnée adimensionnelle
d'action qui intervient ensuite dans la clôture bicouche.

Dans la branche déterminantale, la forme normale de Smith produit
\begin{equation}
 \Sigma_H=\varphi\alpha_{\mathrm{HMT}}^{16},
 \qquad
 \varepsilon_H=\operatorname{ord}_{10}(\Sigma_H)=-34.
 \label{eq:sucesor-102-c14-decada-accion}
\end{equation}
Les deux branches confluent, et seulement alors, dans la grandeur d'action
\begin{equation}
 \boxed{\hbar_{\mathrm{ret}}
 =\eta_{\mathrm{ret}}10^{\varepsilon_H}U_S
 =\eta_{\mathrm{ret}}10^{-34}U_S.}
 \label{eq:sucesor-102-c14-hbar-tipado}
\end{equation}

\section{Publication parangulaire et compatibilité des cartes globale et analytique}
\label{sec:sucesor-102-c14-publicacion-parangular}

Une fois construites la phase \(\alpha_{\mathrm{HMT}}\), la coordonnée
\(\eta_{\mathrm{ret}}\) et la grandeur \(\hbar_{\mathrm{ret}}\), l'application
parangulaire publie
\begin{equation}
 \mathcal P(\alpha_{\mathrm{HMT}},\eta_{\mathrm{ret}})
 =\bigl(A_{\deg},C^*_{\deg}\bigr)
 =\bigl([10^3\alpha_{\mathrm{HMT}}]_{360},
        [2(\eta_{\mathrm{ret}}+\alpha_{\mathrm{HMT}})]_{360}\bigr).
 \label{eq:sucesor-102-c14-mapa-parangular}
\end{equation}
Le facteur \(10^3\) procède de la complétion nonadique ; le facteur \(2\),
de la clôture bidirectionnelle de degré \(2\) ; et \(360\) est la carte globale
qui rend simultanément visibles
\(12\cdot30=4\cdot90=3\cdot120=360\).
La carte analytique applique une seule fois
\begin{equation}
 \kappa_{360}(a,c)=\frac{\pi}{180}(a,c),
 \qquad
 A_{\mathrm{rad}}=\frac{50\pi}{9}\alpha_{\mathrm{HMT}},
 \qquad
 C^*_{\mathrm{rad}}
 =\frac{\pi}{90}(\eta_{\mathrm{ret}}+\alpha_{\mathrm{HMT}}).
 \label{eq:sucesor-102-c14-cartas}
\end{equation}
Par conséquent,
\(C^*_{\deg}=2(\eta_{\mathrm{ret}}+\alpha_{\mathrm{HMT}})\) est le
théorème direct, tandis que
\(\eta_{\mathrm{ret}}=C^*_{\deg}/2-\alpha_{\mathrm{HMT}}\) constitue
exclusivement sa vérification inverse. L'action satisfait
\begin{equation}
 S(\theta)=\hbar_{\mathrm{ret}}\theta_{\mathrm{rad}}
 =h_{\mathrm{ret}}\frac{\theta_{\deg}}{360},
 \qquad h_{\mathrm{ret}}=2\pi\hbar_{\mathrm{ret}}.
 \label{eq:sucesor-102-c14-accion-cartas}
\end{equation}
Les canaux \(q_\pm\) évaluent ultérieurement les cardinaux d'incidence
\(90/120\), déjà générés dans le Topological Phase Kernel ; ils n'en constituent pas
l'origine.
